Pour chacune des affirmations suivantes, indiquer si elle est vraie ou fausse.

Chaque réponse doit être justifiée. Une réponse non justifiée ne rapporte aucun point.

Dans cet exercice, les questions sont indépendantes les unes des autres.

\begin{enumerate}
\item Soit $f$ la fonction définie sur $\R$ par $f(x) = (x + 1)\e^{-x}$.

\textbf{Affirmation 1~:} La limite de la fonction $f$ en $+\infty$ est $+\infty$.

\item On considère la fonction $g$ définie sur $\R$ par $g(x) = x\e^x$.

\textbf{Affirmation 2~:} La fonction $g$ est convexe sur l'intervalle $[-2~;~+\infty[$.

\item \textbf{Affirmation 3~:} Une intégration par parties permet d'obtenir~:

\[\displaystyle\int_{0}^{\pi} x\sin x\,\text{d}x = -\pi.\]

\item On considère l'équation différentielle $(E)$ : $y' = 2y - 5$ et la fonction $h$ définie sur $\R$ par $h(x) = -\dfrac{3}{2}\e^{2x} + \dfrac{5}{2}$.

\textbf{Affirmation 4~:} La fonction $h$ est la solution de $(E)$ qui vaut $1$ en $0$.

\item On considère la fonction \texttt{mystere} définie ci-dessous en Python~:

\begin{center}












\begin{tabular}{|l|}\hline
\texttt{def mystere(n):}\\
\quad \texttt{ f = 1}\\
\quad \texttt{ for k in range(1, n+1):}\\
\quad \quad \texttt{f = f * k}\\
\quad  \texttt{return f}\\ \hline
\end{tabular}

\end{center}

\textbf{Affirmation 5~:} Le nombre renvoyé lors de l'exécution de \texttt{mystere(5)} est $120$.
\end{enumerate}
